\documentclass[10pt,a4paper]{amsart}
\usepackage[margin=19mm]{geometry}
\usepackage{amsmath,amssymb,mathtools}
\usepackage[scr=rsfs,cal=euler]{mathalfa}
\usepackage{xfrac}
\usepackage[unicode,hidelinks,bookmarks=false]{hyperref}
\newcommand{\ie}{i.e.\ }
\newcommand{\cf}{cf.\ }
\newcommand{\resp}{respectively\ }
\newcommand{\Subsection}{\subsection}
\providecommand{\nobreakdash}{\nobreak\hbox{-}}
\setlength{\emergencystretch}{2em}
\multlinegap0pt
\usepackage{bm}
\usepackage{bbm}
\usepackage{dsfont}
\usepackage{xspace}
\usepackage{cancel}
\usepackage{mathtools}
\usepackage{amsthm}
\makeatletter
\@ifundefined{theorem}{\newtheorem{theorem}{Theorem}}{}
\@ifundefined{proposition}{\newtheorem{proposition}[theorem]{Proposition}}{}
\makeatother
\newcommand{\Jac}{\operatorname{Jac}}
\title[The Gromov--Ros conjecture for rank-one symmetric spaces of ]{The Gromov--Ros conjecture for rank-one symmetric spaces of noncompact type}
\author{David Kalaj}
\date{Source: arXiv:2607.20055v3 (CC BY 4.0)}
\begin{document}
\maketitle

\noindent\emph{Source and licence.} The Gromov--Ros conjecture for rank-one symmetric spaces of noncompact type. Version arXiv:2607.20055v3 (CC BY 4.0), distributed under the Creative Commons Attribution 4.0 International licence, as recorded in the arXiv licence metadata for this exact version and shown on the abstract page https://arxiv.org/abs/2607.20055v3. Full rights notice: licenses/arxiv-2607.20055-CC-BY-4.0.txt.\par\smallskip

\noindent\textit{Editorial scope.} Counted unit: the Proposition whose source label is \texttt{prop:basic-matrix} in the article (source0.tex lines 2474--2536, source label \texttt{prop:basic-matrix}), delivered together with the article's own complete original proof of it (source0.tex lines 2538--2752, one \texttt{proof} environment of 215 lines). No mathematics is written, repaired or re-derived: statement and proof are the article's own text line for line. Required context retained uncounted: none. Every definition, display and label of those lines is retained unchanged. Omitted from this package and not redistributed: 1--2473, 2537--2537, 2753--6923. Nothing inside a retained excerpt is omitted or altered: every retained body is delivered exactly as the article's lines are written, so no omission receipt of any kind accompanies this package. Citations: the retained excerpts carry no citation command at all, so no bibliography entry of the article is transcribed and no reference is redistributed. Reference adaptation: none was needed and none was made; every reference inside the delivered unit resolves inside it, so no unresolved reference or citation is produced. Printed slips kept literal: no printed word, symbol, hypothesis or number of the counted statement or of its proof was altered. Layout and class: the article's own class is \texttt{amsart} with packages including fontenc, inputenc, lmodern, microtype, amsmath, amssymb, amsfonts, amsthm, mathtools, mathrsfs, geometry, enumitem, xcolor, hyperref; the wrapper is the lane's pinned \texttt{amsart} editorial wrapper at 10pt on A4, which restates the theorem-like environments the retained text needs and adds the article's own macro definitions that the retained text needs (\texttt{\textbackslash Jac}), with extra wrapper packages (bm, bbm, dsfont, xspace, cancel, mathtools). The delivered numbering is the unit's own. Assets: no retained span contains a figure environment, a graphics-inclusion command, a PDF-page inclusion, a table or a TikZ picture; no image file is copied into this package, the package declares no asset dependency and no image file is redistributed with it. Licence: version arXiv:2607.20055v3 of this article is distributed under the Creative Commons Attribution 4.0 International licence, as recorded in the arXiv licence metadata for this exact version and shown on the abstract page https://arxiv.org/abs/2607.20055v3. A full rights notice is delivered at licenses/arxiv-2607.20055-CC-BY-4.0.txt. Characters: every retained excerpt is pure ASCII, so the delivered engine, XeLaTeX with its default Latin Modern text font, reports no missing character.
\begin{proposition}
\label{prop:basic-matrix}
At every point $(\rho,\theta)$ with $\rho>0$, set
\[
\rho'=u(\theta)\rho,
\qquad
u=e^{t\varphi(\theta)},
\qquad
a=1+\rho^{2},
\qquad
a_t=1+\rho'^2=1+\rho^{2}u^{2}.
\]
Let
\[
N,E_{1},\ldots,E_{n-2},T_{\rho}
\]
and
\[
N',E_{1}',\ldots,E_{n-2}',T_{\rho'}'
\]
denote the source and target orthonormal polar frames, respectively.
Then
\[
DF_{t,\varphi}
=
\begin{pmatrix}
u\sqrt{\dfrac{a}{a_t}}
&
\dfrac{ut}{\sqrt{a_t}}h^{\mathsf T}
&
\dfrac{ut}{\sqrt{aa_t}}q
\\[8pt]
0
&
uI_{n-2}
&
0
\\
0
&
0
&
u\sqrt{\dfrac{a_t}{a}}
\end{pmatrix},
\]
where
\[
h=\nabla_{\mathcal H}\varphi,
\qquad
q=\bar T\varphi.
\]
In particular,
\begin{equation}
\Jac F_{t,\varphi}
=
\det DF_{t,\varphi}
=
u^n
=
e^{nt\varphi}.
\label{eq:basic-ambient-jacobian}
\end{equation}
\end{proposition}

\begin{proof}
The deformation is
\[
F_{t,\varphi}(\rho,\theta)
=
\bigl(\rho',\theta\bigr),
\qquad
\rho'=u(\theta)\rho.
\]
Hence, for a coordinate tangent vector
\[
(s,Y)
\in
\mathbb R\,\partial_\rho
\oplus
T_\theta\mathbb S_o^{n-1},
\]
we have
\[
DF_{t,\varphi}(s,Y)
=
\left(
us+\rho\,du(Y),\,Y
\right).
\]
Since
\[
u=e^{t\varphi},
\qquad
du=ut\,d\varphi,
\]
it follows that
\begin{equation}
DF_{t,\varphi}(s,Y)
=
\left(
us+\rho ut\,d\varphi(Y),\,Y
\right).
\label{eq:master-differential}
\end{equation}

We now express \eqref{eq:master-differential} in the source and
target orthonormal polar frames.

At the source radius $\rho$, the radial unit vector is
\[
N=\sqrt a\,\partial_\rho.
\]
At the target radius $\rho'=u\rho$, it is
\[
N'=\sqrt{a_t}\,\partial_{\rho'}.
\]
Equivalently,
\[
\partial_{\rho'}=\frac{1}{\sqrt{a_t}}N'.
\]

For the radial input $N$, equation
\eqref{eq:master-differential} gives
\[
DF_{t,\varphi}(N)
=
DF_{t,\varphi}\bigl(\sqrt a\,\partial_\rho\bigr)
=
u\sqrt a\,\partial_{\rho'}.
\]
Therefore,
\begin{equation}
DF_{t,\varphi}(N)
=
u\sqrt{\frac{a}{a_t}}\,N'.
\label{eq:DF-radial}
\end{equation}

Next, choose a round-orthonormal horizontal basis
\[
\widehat E_1,\ldots,\widehat E_{n-2}
\in\mathcal H_\theta.
\]
Since the horizontal part of the ambient metric is
\[
\rho^2 g_{\mathcal H},
\]
the corresponding ambient unit vectors at the source are
\[
E_i=\frac{1}{\rho}\widehat E_i.
\]
At the target radius $\rho'=u\rho$, the corresponding ambient unit
vectors are
\[
E_i'=\frac{1}{u\rho}\widehat E_i.
\]

Applying \eqref{eq:master-differential} with
\[
s=0,
\qquad
Y=\frac{1}{\rho}\widehat E_i,
\]
the radial coordinate component is
\[
\rho ut\,
d\varphi\left(\frac{1}{\rho}\widehat E_i\right)
=
ut\,d\varphi(\widehat E_i).
\]
Its coefficient in the target radial unit direction is therefore
\[
\frac{ut}{\sqrt{a_t}}\,
d\varphi(\widehat E_i).
\]
Because $\widehat E_i$ is horizontal,
\[
d\varphi(\widehat E_i)
=
\left\langle
\nabla_{\mathcal H}\varphi,\widehat E_i
\right\rangle
=
\langle h,\widehat E_i\rangle.
\]

The angular component satisfies
\[
\frac{1}{\rho}\widehat E_i
=
u\left(\frac{1}{u\rho}\widehat E_i\right)
=
uE_i'.
\]
Consequently,
\begin{equation}
DF_{t,\varphi}(E_i)
=
\frac{ut}{\sqrt{a_t}}
\langle h,\widehat E_i\rangle N'
+
uE_i'.
\label{eq:DF-horizontal}
\end{equation}

Finally, the source unit Reeb vector is
\[
T_\rho
=
\frac{1}{\rho\sqrt a}\bar T,
\]
whereas at the target radius it is
\[
T_{\rho'}'
=
\frac{1}{u\rho\sqrt{a_t}}\bar T.
\]
Applying \eqref{eq:master-differential} with
\[
s=0,
\qquad
Y=\frac{1}{\rho\sqrt a}\bar T,
\]
and using
\[
q=\bar T\varphi,
\]
the radial coordinate component becomes
\[
\rho ut\,
d\varphi\left(
\frac{1}{\rho\sqrt a}\bar T
\right)
=
\frac{ut}{\sqrt a}\,q.
\]
After converting $\partial_{\rho'}$ into the target radial unit
vector $N'$, this gives
\[
\frac{ut}{\sqrt{aa_t}}q\,N'.
\]

The angular component satisfies
\[
\frac{1}{\rho\sqrt a}\bar T
=
u\sqrt{\frac{a_t}{a}}\,
\frac{1}{u\rho\sqrt{a_t}}\bar T
=
u\sqrt{\frac{a_t}{a}}\,T_{\rho'}'.
\]
Therefore,
\begin{equation}
DF_{t,\varphi}(T_\rho)
=
\frac{ut}{\sqrt{aa_t}}q\,N'
+
u\sqrt{\frac{a_t}{a}}\,T_{\rho'}'.
\label{eq:DF-Reeb}
\end{equation}

Equations \eqref{eq:DF-radial}, \eqref{eq:DF-horizontal}, and
\eqref{eq:DF-Reeb} give the three block columns of the asserted
matrix.

Since this matrix is upper triangular,
\[
\det DF_{t,\varphi}
=
u\sqrt{\frac{a}{a_t}}\,
u^{n-2}\,
u\sqrt{\frac{a_t}{a}}
=
u^n
=
e^{nt\varphi}.
\]
This concludes the proof.
\end{proof}

\end{document}
